\documentclass{article}
\usepackage{amsmath, amssymb, amsthm}
\usepackage{hyperref}
\usepackage{geometry}
\geometry{margin=1in}

\title{Erd\H{o}s Problem \#430}
\date{Last updated: 2025-08-31}
\author{Source: erdosproblems.com}

\begin{document}
\maketitle

\noindent\textbf{Status:} OPEN \quad \textbf{No prize}

\noindent\textbf{Tags:} number theory

\medskip

\noindent\textbf{Problem Statement:}

Fix some integer $n$ and define a decreasing sequence in $[1,n)$ by $a_1=n-1$ and, for $k\geq 2$, letting $a_k$ be the greatest integer in $[1,a_{k-1})$ such that all of the prime factors of $a_k$ are $&#62;n-a_k$. Is it true that, for sufficiently large $n$, not all of this sequence can be prime?




Erd\H{o}s and Graham write 'preliminary calculations made by Selfridge indicate that this is the case but no proof is in sight'. For example if $n=8$ we have $a_1=7$ and $a_2=5$ and then must stop. Sarosh Adenwalla has observed that this problem is equivalent to (the first part of) [385] . Indeed, assuming a positive answer to that, for all large $n$, there exists a composite $m&lt;n$ such that all primes dividing $m$ are $&gt;n-m$. It follows that such an $m$ is equal to some $a_i$ in the sequence defined for $[1,n)$, and $m$ is composite by assumption.

\noindent\textbf{Related OEIS sequences:} possible


\bigskip
\noindent\small{Source: \url{https://www.erdosproblems.com/430}}
\end{document}
